\section{The exact Euler--Navier--Stokes force map and its terminal obstruction}
\label{sec:euler-viscosity-bridge}

All spatial variables in this section remain the original Cartesian variables
$x=(x_1,x_2,x_3)\in\mathbb R^3$.  Time remains $t\in[0,T)$, where
$0<T<\infty$.  The viscosity is an arbitrary fixed real number $\nu>0$;
it is not set equal to one.  Vector norms are Euclidean, matrix norms are
Frobenius, and all $L^p$ norms below are over $\mathbb R^3$.  We use
$(\operatorname{curl}v)_i=\sum_{j,k=1}^3\varepsilon_{ijk}\partial_jv_k$,
$\operatorname{div}v=\sum_j\partial_jv_j$, and
$\Delta v=\sum_j\partial_j^2v$, component by component.  In particular,
$\varepsilon_{123}=1$ fixes the orientation.

\paragraph{Source and the exact scope of the application.}
The released manuscript \emph{Blowup for the Euler equations with smooth
forcing}, Theorem 1.1, states the existence of a smooth Euler solution with
compactly supported smooth initial velocity, a force smooth through a finite
terminal time and supported in a fixed solid torus, finite energy and bounded
spatial velocity gradients on each shorter time interval, and
$\|\operatorname{curl}u(t)\|_\infty\longrightarrow\infty$ as $t\uparrow T$.
Its Section 13.2 and Proposition 13.3 additionally state that the constructed
velocity has fixed compact support and is uniformly bounded before $T$.
The released \texttt{euler-blowup/README.md} states the finite-energy Lipschitz
class and the same vorticity divergence.  These are the specific source facts
used for the application below.  This section proves the force map and the
obstruction to transferring that \emph{same velocity} with a regular force.
It does not repeat or independently certify the manuscript's complete Euler
existence construction or its claimed Lean verification.  No Navier--Stokes
existence or singularity assertion is assumed in the proofs below.

\subsection{The map, including every pressure adjustment}

Let $(u,p_E,f_E)$ solve
\begin{equation}\label{eq:bridge-euler}
 \partial_tu+(u\cdot\nabla)u+\nabla p_E=f_E,
 \qquad \operatorname{div}u=0.
\end{equation}
For any scalar distribution $q$ on $\mathbb R^3\times(0,T)$ define
\begin{equation}\label{eq:bridge-map}
 p_{NS}=p_E+q,\qquad f_{NS}=f_E-\nu\Delta u+\nabla q.
\end{equation}
Then exactly the same velocity $u$ satisfies
\begin{equation}\label{eq:bridge-ns}
 \partial_tu+(u\cdot\nabla)u+\nabla p_{NS}
       =\nu\Delta u+f_{NS},\qquad \operatorname{div}u=0.
\end{equation}
Indeed the left side of \eqref{eq:bridge-ns} is $f_E+\nabla q$ by
\eqref{eq:bridge-euler}, and its right side is
$\nu\Delta u+f_E-\nu\Delta u+\nabla q=f_E+\nabla q$.
Conversely, subtracting \eqref{eq:bridge-euler} from
\eqref{eq:bridge-ns} proves \eqref{eq:bridge-map} with
$q=p_{NS}-p_E$.  Thus \eqref{eq:bridge-map} describes \emph{all} forces
and pressures for which the original velocity solves the viscous equation;
there is no restriction on the pressure adjustment beyond the meaning of
the displayed distributions.

Put $\omega=\operatorname{curl}u$.  Commutation of distributional partial
derivatives gives
\begin{equation}\label{eq:bridge-curl-force}
 \operatorname{curl}f_{NS}
     =\operatorname{curl}f_E-\nu\Delta\omega.
\end{equation}
To verify the pressure cancellation componentwise, write its $i$th component as
\[
(\operatorname{curl}\nabla q)_i
=\sum_{j,k}\varepsilon_{ijk}\partial_j\partial_kq.
\]
the terms indexed by $(j,k)$ and $(k,j)$ cancel because the derivatives
commute and $\varepsilon_{ikj}=-\varepsilon_{ijk}$.
The identity \eqref{eq:bridge-curl-force} is therefore independent of $q$,
including pressure adjustments whose gradients have bad behavior at
spatial infinity.  An arbitrary distribution $q$ need not possess time
sections.  Nevertheless, since $u$ and $f_E$ are smooth before $T$, the
right side of \eqref{eq:bridge-curl-force} is smooth and is the canonical
representative of the distribution $\operatorname{curl}f_{NS}$.  Its
time sections are therefore defined for every $t<T$, independently of
whether $q$ or $f_{NS}$ has time sections.

\subsection{A terminal energy bound from the stated preterminal class}

Here is a complete justification that the source's preterminal energy
hypothesis supplies a bound independent of the shorter interval.  Suppose
$u$ is smooth on $\mathbb R^3\times[0,T)$ and solves
\eqref{eq:bridge-euler}, with a distributional pressure, and for every
$\tau<T$ satisfies
\begin{equation}\label{eq:bridge-energy-class}
 u\in L^\infty([0,\tau];L^2),\qquad
 \nabla u\in L^\infty(\mathbb R^3\times[0,\tau]).
\end{equation}
Suppose $f_E$ is continuous into $L^2$ on every such shorter interval and
$\int_0^T\|f_E(s)\|_2\,ds<\infty$.  The compactly supported smooth
force in the source has these properties: a common compact containing set
$K$ gives
$\|f_E(s)\|_2\le |K|^{1/2}\sup_{K\times[0,T]}|f_E|$.

First, a spatially Lipschitz field $v\in L^2$ with Lipschitz constant
$L=\|\nabla v\|_\infty>0$ is bounded.  If $A=|v(x)|>0$, then
$|v(y)|\ge A/2$ on the ball of radius $A/(2L)$ about $x$.  Hence
\[
 \|v\|_2^2\ge \frac{A^2}{4}\frac{4\pi}{3}
                    \left(\frac{A}{2L}\right)^3
       =\frac{\pi A^5}{24L^3},\qquad
 \|v\|_\infty\le
       \left(\frac{24}{\pi}\right)^{1/5}\|v\|_2^{2/5}L^{3/5}.
\]
If $L=0$, $v$ is constant and its finite $L^2$ norm forces $v=0$.
Smoothness and Fatou's lemma extend the almost-every-time bounds in
\eqref{eq:bridge-energy-class} to every time section.  The inequality
therefore bounds $u$ uniformly on each $[0,\tau]$.

We justify the elimination of pressure without a growth assumption on
$p_E$.  Use the unitary Fourier convention
$\widehat v(\xi)=(2\pi)^{-3/2}\int e^{-ix\cdot\xi}v(x)\,dx$,
extended by density to $L^2$.  Let $P$ be the orthogonal $L^2$ projection
with multiplier
$I-\xi\otimes\xi/|\xi|^2$ for $\xi\ne0$; its value at $0$ does not
affect an $L^2$ function.  This real symmetric idempotent multiplier has
operator norm at most one.  It projects onto the fields with
distributional divergence zero and preserves curl, as is verified by
$\xi\cdot Pv=0$ and $\xi\times Pv=\xi\times v$.

Set $g=f_E-(u\cdot\nabla)u$.  On $[0,\tau]$ this is an $L^\infty_tL^2_x$
field, since
$\|(u\cdot\nabla)u\|_2\le\|\nabla u\|_\infty\|u\|_2$.
Taking curl of Euler and using $\operatorname{div}u=0$ shows that
both curl and divergence of the space-time distribution
$\partial_tu-Pg$ vanish.  For a compactly supported smooth time test
$\eta$ in $(0,\tau)$, its pairing with that test is the $L^2$ field
\[
 V_\eta=-\int_0^\tau u(t)\eta'(t)\,dt
             -\int_0^\tau Pg(t)\eta(t)\,dt.
\]
It has zero curl and divergence.  Fourier transformation gives
$\xi\times\widehat V_\eta=0$ and
$\xi\cdot\widehat V_\eta=0$ almost everywhere.  For $\xi\ne0$ the
first relation makes $\widehat V_\eta$ parallel to $\xi$ and the second
makes it perpendicular, so it is zero.  The remaining point $\xi=0$
has measure zero.  Thus $V_\eta=0$, proving
$\partial_tu=Pg$ as an $L^2$-valued distribution in time.

For completeness, the distributional conclusion supplies an absolutely
continuous $L^2$ representative: subtract the Bochner primitive
$\int_0^t Pg(s)\,ds$ from $u$; pairing the remainder against each
$L^2$ test vector gives a scalar distribution with derivative zero and
hence a constant.  The constants define one $L^2$ vector, using the
uniform $L^2$ bounds.  This gives the claimed representative.  Its
strong $L^2$ continuity and the original local smoothness identify it
with the given $u(t)$ at every $t$, including $t=0$: choose times of
almost-everywhere agreement converging to the chosen time and pair
against compactly supported spatial tests.  Consequently
$u(t)=u(0)+\int_0^t Pg(s)\,ds$ in $L^2$.

The Hilbert-space norm identity now gives, for almost every $t<\tau$,
\[
 \frac12\frac{d}{dt}\|u(t)\|_2^2
    =\langle u,Pg\rangle_{L^2}
    =\langle u,f_E\rangle_{L^2}
      -\int u\cdot((u\cdot\nabla)u)\,dx.
\]
The norm identity follows by expanding the square of
$u(t+h)=u(t)+\int_t^{t+h}Pg(s)\,ds$ and using the Lebesgue
differentiation theorem for the $L^2$-valued integrand.  Orthogonality
of $P$ and $Pu=u$ justify its removal from the inner product.

Choose a smooth cutoff $\chi_R(x)=\chi(x/R)$, equal to one on
$|x|\le R$, zero for $|x|\ge2R$, with $0\le\chi\le1$ and
$|\nabla\chi_R|\le C_\chi/R$.  Local integration by parts and
$\operatorname{div}u=0$ give
\[
 \int\chi_R u\cdot((u\cdot\nabla)u)\,dx
       =-\frac12\int|u|^2u\cdot\nabla\chi_R\,dx.
\]
The right side has absolute value at most
$C_\chi\|u\|_\infty\|u\|_2^2/(2R)$ and tends to zero.
The left integrand is dominated by
$\|\nabla u\|_\infty|u|^2$, so dominated convergence applies.
We have proved the exact Euler energy identity
\begin{equation}\label{eq:bridge-energy-identity}
 \frac12\|u(t)\|_2^2-\frac12\|u(0)\|_2^2
          =\int_0^t\langle u(s),f_E(s)\rangle_{L^2}\,ds.
\end{equation}
For $\epsilon>0$, differentiation of
$(\|u(t)\|_2^2+\epsilon^2)^{1/2}$, the identity just proved, and
Cauchy--Schwarz bound its derivative above by $\|f_E(t)\|_2$.
Integrating and then taking $\epsilon\downarrow0$ proves
\begin{equation}\label{eq:bridge-E}
 \|u(t)\|_2\le\|u(0)\|_2+\int_0^t\|f_E(s)\|_2\,ds
 \le E:=\|u(0)\|_2+\int_0^T\|f_E(s)\|_2\,ds<\infty.
\end{equation}
The proof works on every $[0,\tau]$, and its final constant does not
depend on $\tau$.  It therefore proves \eqref{eq:bridge-E} for every
$t<T$.  For the actual constructed field, the fixed support and
uniform velocity bound in Proposition 13.3 give a second direct
verification of this conclusion.

\subsection{A heat identity with no hidden assumption at spatial infinity}

Let $v\in L^2(\mathbb R^3;\mathbb R^3)$ and put
$\zeta=\operatorname{curl}v$ as a tempered distribution.  Suppose the
distribution $\Delta\zeta$ is represented by a bounded vector field $h$.
Define, for $a>0$,
\[
 G_a(x)=(4\pi a)^{-3/2}\exp\left(-\frac{|x|^2}{4a}\right),
 \qquad e^{a\Delta}w=G_a*w.
\]
The following identity is an equality of tempered distributions, with a
bounded continuous representative on its right:
\begin{equation}\label{eq:bridge-heat-identity}
 \zeta=\operatorname{curl}(G_a*v)
                 -\int_0^a G_r*h\,dr.
\end{equation}
In particular,
\begin{equation}\label{eq:bridge-heat-bound}
 \|\operatorname{curl}v\|_\infty
 \le C_0a^{-5/4}\|v\|_2+a\|\Delta\operatorname{curl}v\|_\infty,
 \qquad C_0=\frac{1}{\sqrt2}(8\pi)^{-3/4},\quad a>0.
\end{equation}
For a smooth $v$ the bound concerns its ordinary pointwise curl.

Here are all convergence and constant computations.  A function in $L^2$
defines a tempered distribution by Cauchy--Schwarz on Schwartz test
functions; its distributional derivatives do so as well.  For a Schwartz
test function $\phi$, the map $r\mapsto G_r*\phi$, with value $\phi$
at $r=0$, is continuously differentiable in Schwartz space on every
bounded interval of nonnegative $r$, and its derivative is
$G_r*\Delta\phi$.  This follows directly by Fourier transformation:
the multipliers are $e^{-r|\xi|^2}$ and
$-|\xi|^2e^{-r|\xi|^2}$; differentiating each Schwartz seminorm is
justified by the rapid decay of $\widehat\phi$ and polynomial bounds
on all derivatives of these multipliers, uniform for bounded $r$.
The fundamental theorem of calculus, paired with $\zeta$, gives
\[
 \langle G_a*\zeta-\zeta,\phi\rangle
       =\int_0^a\langle\Delta\zeta,G_r*\phi\rangle\,dr
       =\int_0^a\langle G_r*h,\phi\rangle\,dr.
\]
Fubini is valid because
$|\langle h,G_r*\phi\rangle|
 \le\|h\|_\infty\|\phi\|_1$ for every $r\ge0$;
the Gaussian has integral one and is nonnegative.
Differentiating the convolution with $v$ in distributions gives
$G_a*\zeta=\operatorname{curl}(G_a*v)$, proving
\eqref{eq:bridge-heat-identity}.  The function
$\int_0^aG_r*h\,dr$ is bounded by $a\|h\|_\infty$.
It is continuous: integrals over $[\delta,a]$ are continuous, and the
integral over $[0,\delta]$ has uniform norm at most
$\delta\|h\|_\infty$, which tends to zero.  This proves the asserted
representative without a pointwise or growth assumption on $\zeta$.

For a real unit vector $b$ and every $x$, the scalar triple product
and Cauchy--Schwarz give
\[
 \left|b\cdot\operatorname{curl}(G_a*v)(x)\right|
   =\left|\int v(y)\cdot\bigl(b\times\nabla G_a(x-y)\bigr)\,dy\right|
   \le\|v\|_2\left(\int|b\times\nabla G_a(z)|^2\,dz\right)^{1/2}.
\]
The sign in this formula follows from
$\operatorname{curl}(G_a*v)=\int\nabla G_a(x-y)\times v(y)\,dy$.
The Gaussian integrals are
\[
 \|G_a\|_2^2
   =(4\pi a)^{-3}(2\pi a)^{3/2}=(8\pi a)^{-3/2},
 \qquad
 \int(\partial_jG_a)(\partial_kG_a)
       =\frac{\delta_{jk}}{4a}\|G_a\|_2^2.
\]
Indeed $\partial_jG_a=-x_jG_a/(2a)$, odd integrands have zero
integral, and integration by parts in one dimension gives
$\int x_j^2G_a^2=a\int G_a^2$.
Thus
\[
 \int|b\times\nabla G_a|^2
  =\int\bigl(|\nabla G_a|^2-|b\cdot\nabla G_a|^2\bigr)
  =\frac{1}{2a}(8\pi a)^{-3/2}.
\]
Taking the supremum over $|b|=1$ yields
$\|\operatorname{curl}(G_a*v)\|_\infty
 \le C_0a^{-5/4}\|v\|_2$ with precisely the constant displayed in
\eqref{eq:bridge-heat-bound}.  Together with
\eqref{eq:bridge-heat-identity} this proves the bound.

There is no surviving spatial harmonic term in this argument: the
distributional identity started with $\zeta=\operatorname{curl}v$ for
$v\in L^2$, kept its exact heat evolution, and integrated its actual
Laplacian.  No boundedness of $\zeta$ was used to prove its boundedness.

\subsection{The pressure-independent obstruction and its rate}

Apply \eqref{eq:bridge-heat-bound} to each original time section $u(t)$.
Write
\[
 M(t)=\|\omega(t)\|_\infty,\quad
 F_E(t)=\|\operatorname{curl}f_E(t)\|_\infty,\quad
 F_{NS}(t)=\|(\operatorname{curl}f_E-\nu\Delta\omega)(t)\|_\infty.
\]
Thus $F_{NS}$ always uses the canonical curl representative just proved
to exist, even for an arbitrary distributional pressure adjustment.
Whenever $F_{NS}(t)<\infty$, identity
\eqref{eq:bridge-curl-force} shows that
$\Delta\omega(t)$ is bounded and that
\begin{equation}\label{eq:bridge-direct-bound}
 M(t)\le C_0a^{-5/4}E+
                    \frac{a}{\nu}\bigl(F_E(t)+F_{NS}(t)\bigr)
 \qquad\hbox{for every }a>0.
\end{equation}
In particular, a finite uniform bound on the two force curls gives a
finite uniform bound on vorticity, contradicting the source's
$M(t)\longrightarrow\infty$.

The conclusion already excludes uniformly bounded first spatial
derivatives of $f_{NS}$.  Explicitly, for any vector field $f$,
\[
 |\operatorname{curl}f|^2
   =(\partial_2f_3-\partial_3f_2)^2
    +(\partial_3f_1-\partial_1f_3)^2
    +(\partial_1f_2-\partial_2f_1)^2
   \le2|\nabla f|_F^2,
\]
using $(a-b)^2\le2(a^2+b^2)$ on each displayed summand.
The source's fixed support and smoothness through $T$ give
$F_E^*:=\sup_{t<T}F_E(t)<\infty$.
If $f_{NS}$ is a smooth spacetime force and
$\sup_{t<T}\|\nabla f_{NS}(t)\|_\infty=B<\infty$, then
\eqref{eq:bridge-direct-bound} gives
\[
 M(t)\le C_0a^{-5/4}E+
                  \frac{a}{\nu}\bigl(F_E^*+\sqrt2 B\bigr)
 \quad(0\le t<T)
\]
for each fixed $a>0$, an explicit contradiction.
More generally, suppose the spacetime distributions
$\partial_j(f_{NS})_k$ have bounded measurable representatives, whose
matrix has essential supremum Frobenius norm $B$ on
$\mathbb R^3\times(0,T)$.  The same componentwise inequality gives a
spacetime essential bound $\sqrt2B$ on the curl.  The canonical curl
representative is smooth, so this bound holds everywhere: a strict
violation at one point would persist on an open set of positive
measure.  The displayed contradiction therefore also rules out
uniformly bounded spacetime weak first derivatives for arbitrary
distributional forces, without assuming time sections of their gradients.

One can also quantify the necessary loss.  The source has $E>0$:
otherwise \eqref{eq:bridge-E} gives $u=0$ and $M=0$.
Put $A=C_0E>0$.  At a time with $M(t)>0$, rearrange
\eqref{eq:bridge-heat-bound}, using $\|u(t)\|_2\le E$, to get
\[
 \|\Delta\omega(t)\|_\infty
       \ge \frac{M(t)}{a}-\frac{A}{a^{9/4}}.
\]
Take the admissible positive value
$a=(9A/(4M(t)))^{4/5}$.  This is the maximizer: differentiating the
right side gives $-M(t)a^{-2}+(9A/4)a^{-13/4}$, with a unique zero
there, positive derivative before that zero and negative derivative
after it.  Substitution gives
\begin{equation}\label{eq:bridge-laplacian-lower}
 \|\Delta\omega(t)\|_\infty
 \ge \frac{5}{9}\left(\frac{4}{9}\right)^{4/5}
               (C_0E)^{-4/5}M(t)^{9/5}.
\end{equation}
The reverse triangle inequality in \eqref{eq:bridge-curl-force}
therefore proves, for every pressure adjustment,
\begin{equation}\label{eq:bridge-force-lower}
 F_{NS}(t)\ge
 \nu\frac{5}{9}\left(\frac{4}{9}\right)^{4/5}
                (C_0E)^{-4/5}M(t)^{9/5}-F_E(t).
\end{equation}
When a force curl is unbounded at the chosen time the inequality holds
in the extended sense.  If it is bounded, all preceding steps apply.
For $M(t)=0$ the same displayed lower bound is simply
$F_{NS}(t)\ge-F_E(t)$ and also holds.  Finally, whenever the force has
smooth spacetime representatives, as every candidate in the target problem
must have,
\begin{equation}\label{eq:bridge-gradient-lower}
 \|\nabla f_{NS}(t)\|_\infty\ge\frac1{\sqrt2}
 \left[\nu\frac{5}{9}\left(\frac{4}{9}\right)^{4/5}
              (C_0E)^{-4/5}M(t)^{9/5}-F_E(t)\right].
\end{equation}
Consequently $M(t)\to\infty$ and $F_E^*<\infty$ imply
$F_{NS}(t)\to\infty$ for every distributional pressure adjustment in
\eqref{eq:bridge-map}.  For smooth spacetime forces they also imply
$\|\nabla f_{NS}(t)\|_\infty\to\infty$.  This is stronger than
excluding only the choice
$p_{NS}=p_E$.

\paragraph{The precise result for the released construction.}
The original Euler velocity has a well-defined exact viscous realization
with force $f_E-\nu\Delta u+\nabla q$ for every $\nu>0$ and every
pressure adjustment $q$.  Every such force has a canonical smooth curl
whose spatial supremum diverges as the original terminal time is
approached, and none has uniformly bounded spacetime weak first
derivatives.  For smooth spacetime forces,
their spatial first-derivative supremum also diverges at that time.
Thus this realization cannot have a force smooth through the terminal
time with a common compact support, nor can it satisfy a requirement
of globally bounded first spatial derivatives through that time.
On $\mathbb R^3$, smoothness alone on a noncompact spatial domain would
not imply a global derivative bound; the fixed-support or stated decay
requirement is material and has been kept explicit.

This obstruction concerns the original velocity and the exact map
\eqref{eq:bridge-map}.  It proves no nonexistence statement about a
different Navier--Stokes velocity, a different construction, a
viscosity-dependent change of the evolving fields, or a different
mathematical correspondence.  It also gives no Navier--Stokes
counterexample with the initial data and forcing regularity required
by a Millennium problem formulation.  The map and the obstruction,
rather than a mere comparison of equation names, identify exactly
what a successful further transfer must change.
